\begin{afterword}
\summaryhead

The post argues that people can make almost any goal seem to imply almost any action, and that this is hard to catch because we rarely know exactly what a stated criterion recommends. Whether Robert Mugabe would resign if he cared only for Zimbabwe cannot be proved; motives and the world are too complicated. What is needed is an optimizer with a known goal and no other desires, whose output can be compared with human rationalizations of that goal.

Natural selection is such an optimizer: it has no mercy, politics or ideology, and its output is optimized only for inclusive genetic fitness. The group-selection story in the previous post gives the comparison. According to the post, Wynne-Edwards expected group selection to produce restraint in breeding and enough food for all; the laboratory produced cannibalism. The author answers the objection that human morality does not care about fitness: if we cannot see what a single simple criterion implies, we cannot judge complex ones, so we need simple practice cases.

\responsehead

The question the post asks is a real one. It is hard to know what a goal implies, and a case where the goal is known would help to measure how far people's reasoning bends. The post also concedes, in a parenthesis, that rationalization can be shown more simply: people's views on office politics track their interests. So its case for studying evolution rests on the second need, a criterion whose implications are known exactly.

The comparison with Wynne-Edwards is loose. The post says he expected ``Voluntary individual restraint in breeding, and enough food for everyone'' from group selection for small population size. I found no such prediction. His hypothesis was about social conventions keeping numbers below what the food will bear, with surplus animals dying, and the experiment came fourteen years after his book. On Wade's results, see the notes on ``The Tragedy of Group Selectionism.''

The lesson is argued by analogy: learn to hear one note before an orchestra. The post gives no evidence that practice on evolutionary cases improves judgement of moral principles. In its own example the answer came from an experiment, not from seeing clearly what the criterion implied, and for human criteria, as the post itself says, there is no such check. Finally, the post explains our talent for rationalizing by ancestral politics, stated as settled (``There's no mystery about this'') with no evidence.

In short: a real question about what a goal implies, answered with a prediction attributed to Wynne-Edwards that could not be found, and a lesson that rests on an analogy.
\end{afterword}
